\documentclass[10pt,a4paper]{amsart}
\usepackage{amsmath,amssymb,amsthm,graphicx}
% Retained preamble for the Stacks Project chapter morphisms.tex (Morphisms of Schemes).
% The chapter file pinned-original-source.tex carries only \input{preamble} on its first line, so no
% preamble text lives inside the chapter itself.  The Stacks Project uses one single shared preamble for
% every chapter; the snapshot stages it verbatim in sources/stacks-chow/preamble.tex.  The block below is
% copied mechanically, line for line, from that file's lines 14-18 (the Xy-pic setup) and lines 154-217
% (the native theorem declarations together with the mathematical macros).  Nothing was added, removed or
% redefined.  The wrapper class and the amsmath/amssymb/amsthm/graphicx base packages are supplied by the
% assembler, so the class/\externaldocument part of the shared preamble is not repeated here.
% For commutative diagrams we use Xy-pic
\usepackage[all]{xy}

% We use 2cell for 2-commutative diagrams.
\xyoption{2cell}
% Theorem environments.
%
\theoremstyle{plain}
\newtheorem{theorem}[subsection]{Theorem}
\newtheorem{proposition}[subsection]{Proposition}
\newtheorem{lemma}[subsection]{Lemma}

\theoremstyle{definition}
\newtheorem{definition}[subsection]{Definition}
\newtheorem{example}[subsection]{Example}
\newtheorem{exercise}[subsection]{Exercise}
\newtheorem{situation}[subsection]{Situation}

\theoremstyle{remark}
\newtheorem{remark}[subsection]{Remark}
\newtheorem{remarks}[subsection]{Remarks}

\numberwithin{equation}{subsection}

% Macros
%
\def\lim{\mathop{\mathrm{lim}}\nolimits}
\def\colim{\mathop{\mathrm{colim}}\nolimits}
\def\Spec{\mathop{\mathrm{Spec}}}
\def\Hom{\mathop{\mathrm{Hom}}\nolimits}
\def\Ext{\mathop{\mathrm{Ext}}\nolimits}
\def\SheafHom{\mathop{\mathcal{H}\!\mathit{om}}\nolimits}
\def\SheafExt{\mathop{\mathcal{E}\!\mathit{xt}}\nolimits}
\def\Sch{\mathit{Sch}}
\def\Mor{\mathop{\mathrm{Mor}}\nolimits}
\def\Ob{\mathop{\mathrm{Ob}}\nolimits}
\def\Sh{\mathop{\mathit{Sh}}\nolimits}
\def\NL{\mathop{N\!L}\nolimits}
\def\CH{\mathop{\mathrm{CH}}\nolimits}
\def\proetale{{pro\text{-}\acute{e}tale}}
\def\etale{{\acute{e}tale}}
\def\QCoh{\mathit{QCoh}}
\def\Ker{\mathop{\mathrm{Ker}}}
\def\Im{\mathop{\mathrm{Im}}}
\def\Coker{\mathop{\mathrm{Coker}}}
\def\Coim{\mathop{\mathrm{Coim}}}

% Boxtimes
%
\DeclareMathSymbol{\boxtimes}{\mathbin}{AMSa}{"02}

%
% Macros for moduli stacks/spaces
%
\def\QCohstack{\mathcal{QC}\!\mathit{oh}}
\def\Cohstack{\mathcal{C}\!\mathit{oh}}
\def\Spacesstack{\mathcal{S}\!\mathit{paces}}
\def\Quotfunctor{\mathrm{Quot}}
\def\Hilbfunctor{\mathrm{Hilb}}
\def\Curvesstack{\mathcal{C}\!\mathit{urves}}
\def\Polarizedstack{\mathcal{P}\!\mathit{olarized}}
\def\Complexesstack{\mathcal{C}\!\mathit{omplexes}}
% \Pic is the operator that assigns to X its picard group, usage \Pic(X)
% \Picardstack_{X/B} denotes the Picard stack of X over B
% \Picardfunctor_{X/B} denotes the Picard functor of X over B
\def\Pic{\mathop{\mathrm{Pic}}\nolimits}
\def\Picardstack{\mathcal{P}\!\mathit{ic}}
\def\Picardfunctor{\mathrm{Pic}}
\def\Deformationcategory{\mathcal{D}\!\mathit{ef}}

% --- appended 2026-10-12 (Chief): the project own definitions of the reference/slogan/history
% display environments, copied mechanically from the shared Stacks preamble lines 8-12
% (\usepackage{verbatim} plus the three \newenvironment lines, which the chapter preamble keeps
% verbatim from the project preamble). Nothing added, removed or redefined: these environments
% carry no mathematical content and the project itself typesets them as comments. Without these
% three lines a counted statement whose own text uses them fails with Environment slogan undefined.
\usepackage{verbatim}
\newenvironment{reference}{\comment}{\endcomment}
\newenvironment{slogan}{\comment}{\endcomment}
\newenvironment{history}{\comment}{\endcomment}
\usepackage{cleveref}
\crefname{originallemma}{original lemma}{original lemmas}
\Crefname{originallemma}{Original lemma}{Original lemmas}
\makeatletter
\newlabel{algebra-lemma-NAK}{{00DV}{source}{}{sourcepointer-1}{}}
\newlabel{algebra-lemma-NAK@cref}{{[originallemma][0][]00DV}{[1][0][]source}{}{}{}}
\makeatother
\crefname{originallemma}{original lemma}{original lemmas}
\Crefname{originallemma}{Original lemma}{Original lemmas}
\makeatletter
\newlabel{algebra-lemma-support-closed}{{00L2}{source}{}{sourcepointer-2}{}}
\newlabel{algebra-lemma-support-closed@cref}{{[originallemma][0][]00L2}{[1][0][]source}{}{}{}}
\makeatother
\crefname{originallemma}{original lemma}{original lemmas}
\Crefname{originallemma}{Original lemma}{Original lemmas}
\makeatletter
\newlabel{lemma-support-affine-open}{{056I}{source}{}{sourcepointer-3}{}}
\newlabel{lemma-support-affine-open@cref}{{[originallemma][0][]056I}{[1][0][]source}{}{}{}}
\makeatother
\crefname{originallemma}{original lemma}{original lemmas}
\Crefname{originallemma}{Original lemma}{Original lemmas}
\makeatletter
\newlabel{modules-lemma-pullback-finite-type}{{01B6}{source}{}{sourcepointer-4}{}}
\newlabel{modules-lemma-pullback-finite-type@cref}{{[originallemma][0][]01B6}{[1][0][]source}{}{}{}}
\makeatother
\crefname{originallemma}{original lemma}{original lemmas}
\Crefname{originallemma}{Original lemma}{Original lemmas}
\makeatletter
\newlabel{modules-lemma-support-finite-type-closed}{{01BA}{source}{}{sourcepointer-5}{}}
\newlabel{modules-lemma-support-finite-type-closed@cref}{{[originallemma][0][]01BA}{[1][0][]source}{}{}{}}
\makeatother
\crefname{originallemma}{original lemma}{original lemmas}
\Crefname{originallemma}{Original lemma}{Original lemmas}
\makeatletter
\newlabel{properties-lemma-finite-type-module}{{01PB}{source}{}{sourcepointer-6}{}}
\newlabel{properties-lemma-finite-type-module@cref}{{[originallemma][0][]01PB}{[1][0][]source}{}{}{}}
\makeatother
\crefname{originallemma}{original lemma}{original lemmas}
\Crefname{originallemma}{Original lemma}{Original lemmas}
\makeatletter
\newlabel{sheaves-lemma-stalk-pullback-modules}{{0098}{source}{}{sourcepointer-7}{}}
\newlabel{sheaves-lemma-stalk-pullback-modules@cref}{{[originallemma][0][]0098}{[1][0][]source}{}{}{}}
\makeatother
\title{The support of a finite type quasi-coherent module is closed, detected on residue fibres, and compatible with pullback}
\author{The Stacks Project Authors}
\date{stacks-project:morphisms.tex:local-snapshot-20261012; GNU FDL 1.2 or later}
\begin{document}
\maketitle
\noindent\textit{Editorial note.} All delivered mathematical text is the author's own, copied from the pinned Stacks Project chapter file pinned-original-source.tex; no statement or proof step was rewritten, reordered or removed.\par\medskip
\noindent\textit{Editorial note.} Following the multi-claim convention for this source, source\_url is the source-level URL https://stacks.math.columbia.edu/ and the source label lemma-support-finite-type is carried in claim\_key and here rather than in a per-claim tag URL.\par\medskip
\noindent\textit{Editorial note.} The preamble is the chapter's shared Stacks Project preamble, retained as sources/stacks-morphisms/retained-preamble.tex: a line-for-line excerpt of the shared Stacks preamble staged with this snapshot. No macro was invented or redefined, and this card adds no preamble text of its own.\par\medskip
\noindent\textit{Editorial note.} Statement span 678-696 and proof span 698-726 come from the halfA pool card operations/pool-cards/stacks-morphisms-halfA.json and were re-read in the pinned chapter file before this card was written.\par\medskip
\noindent\textit{Editorial note.} The proof refers to the three parts of the statement by number; those references are the author's and are preserved.\par\medskip
\noindent\textit{Editorial note.} 7 cross-references occur inside the counted spans; they are listed in external\_references as navigation pointers only.\par\medskip
\noindent\textit{Allowed standard prerequisites (editorial).}\begin{itemize}
\item The support of a finite type quasi-coherent module is closed, cited from the modules chapter, and the affine description of a finite type quasi-coherent module as the tilde of a finite module together with the closedness of supports over affine schemes. Cited at proof lines 915-920. Gives the closedness of the support, in the cited form or by combining the affine statements.
\item Nakayama's lemma in the form that a stalk is zero exactly when its tensor product with the residue field is zero. Cited at proof lines 925-927. Supplies the second equivalence of part (2) and is used again at the end of the proof for the pullback.
\item The stalk of a pullback module is the tensor product of the original stalk with the local ring, cited from the sheaves chapter, and the pullback of a finite type module is of finite type, cited from the modules chapter. Cited at proof lines 932-937. Reduces the compatibility of supports with pullback to the residue fibre criterion of part (2).
\end{itemize}\medskip
\ifdefined\hypertarget\hypertarget{sourcepointer-1}{}\fi
\noindent\textit{Source reference.} Original lemma 00DV: focused-seeds/AGP/sources/stacks-snapshot-20261011/tags.txt, source lines 512--512. Same-chapter reference used in the retained author text at source line(s) 709, 5128, 7341, 7915, 11427, 15408. It is supplied only as a navigation pointer so that the retained text compiles; the counted item delivers the author's own statement and the author's complete proof, and no part of the referenced item is copied in or substituted for the counted argument.\par\medskip
\ifdefined\hypertarget\hypertarget{sourcepointer-2}{}\fi
\noindent\textit{Source reference.} Original lemma 00L2: focused-seeds/AGP/sources/stacks-snapshot-20261011/tags.txt, source lines 772--772. Same-chapter reference used in the retained author text at source line(s) 705, 759, 11631, 11640, 11679, 11688. It is supplied only as a navigation pointer so that the retained text compiles; the counted item delivers the author's own statement and the author's complete proof, and no part of the referenced item is copied in or substituted for the counted argument.\par\medskip
\ifdefined\hypertarget\hypertarget{sourcepointer-3}{}\fi
\noindent\textit{Source reference.} Original lemma 056I: focused-seeds/AGP/sources/stacks-snapshot-20261011/tags.txt, source lines 6711--6711. Same-chapter reference used in the retained author text at source line(s) 702. It is supplied only as a navigation pointer so that the retained text compiles; the counted item delivers the author's own statement and the author's complete proof, and no part of the referenced item is copied in or substituted for the counted argument.\par\medskip
\ifdefined\hypertarget\hypertarget{sourcepointer-4}{}\fi
\noindent\textit{Source reference.} Original lemma 01B6: focused-seeds/AGP/sources/stacks-snapshot-20261011/tags.txt, source lines 1715--1715. Same-chapter reference used in the retained author text at source line(s) 712, 10678. It is supplied only as a navigation pointer so that the retained text compiles; the counted item delivers the author's own statement and the author's complete proof, and no part of the referenced item is copied in or substituted for the counted argument.\par\medskip
\ifdefined\hypertarget\hypertarget{sourcepointer-5}{}\fi
\noindent\textit{Source reference.} Original lemma 01BA: focused-seeds/AGP/sources/stacks-snapshot-20261011/tags.txt, source lines 1719--1719. Same-chapter reference used in the retained author text at source line(s) 700, 5097, 7918. It is supplied only as a navigation pointer so that the retained text compiles; the counted item delivers the author's own statement and the author's complete proof, and no part of the referenced item is copied in or substituted for the counted argument.\par\medskip
\ifdefined\hypertarget\hypertarget{sourcepointer-6}{}\fi
\noindent\textit{Source reference.} Original lemma 01PB: focused-seeds/AGP/sources/stacks-snapshot-20261011/tags.txt, source lines 2224--2224. Same-chapter reference used in the retained author text at source line(s) 703, 754, 4224, 5201, 5286, 6703, 10562, 11247. It is supplied only as a navigation pointer so that the retained text compiles; the counted item delivers the author's own statement and the author's complete proof, and no part of the referenced item is copied in or substituted for the counted argument.\par\medskip
\ifdefined\hypertarget\hypertarget{sourcepointer-7}{}\fi
\noindent\textit{Source reference.} Original lemma 0098: focused-seeds/AGP/sources/stacks-snapshot-20261011/tags.txt, source lines 344--344. Same-chapter reference used in the retained author text at source line(s) 720, 4907. It is supplied only as a navigation pointer so that the retained text compiles; the counted item delivers the author's own statement and the author's complete proof, and no part of the referenced item is copied in or substituted for the counted argument.\par\medskip
\begin{lemma}
\label{lemma-support-finite-type}
Let $\mathcal{F}$ be a finite type quasi-coherent module
on a scheme $X$. Then
\begin{enumerate}
\item The support of $\mathcal{F}$ is closed.
\item For $x \in X$ we have
$$
x \in \text{Supp}(\mathcal{F})
\Leftrightarrow
\mathcal{F}_x \not = 0
\Leftrightarrow
\mathcal{F}_x \otimes_{\mathcal{O}_{X, x}} \kappa(x) \not = 0.
$$
\item For any morphism of schemes $f : Y \to X$ the pullback
$f^*\mathcal{F}$ is of finite type as well and we have
$\text{Supp}(f^*\mathcal{F}) = f^{-1}(\text{Supp}(\mathcal{F}))$.
\end{enumerate}
\end{lemma}

\begin{proof}
Part (1) is a reformulation of
Modules, Lemma \ref{modules-lemma-support-finite-type-closed}.
You can also combine
Lemma \ref{lemma-support-affine-open},
Properties, Lemma \ref{properties-lemma-finite-type-module},
and
Algebra, Lemma \ref{algebra-lemma-support-closed}
to see this. The first equivalence in (2) is the definition
of support, and the second equivalence follows from
Nakayama's lemma, see
Algebra, Lemma \ref{algebra-lemma-NAK}.
Let $f : Y \to X$ be a morphism of schemes. Note that
$f^*\mathcal{F}$ is of finite type by
Modules, Lemma \ref{modules-lemma-pullback-finite-type}.
For the final assertion, let $y \in Y$ with image $x \in X$.
Recall that
$$
(f^*\mathcal{F})_y =
\mathcal{F}_x \otimes_{\mathcal{O}_{X, x}} \mathcal{O}_{Y, y},
$$
see
Sheaves, Lemma \ref{sheaves-lemma-stalk-pullback-modules}.
Hence $(f^*\mathcal{F})_y \otimes \kappa(y)$ is nonzero
if and only if $\mathcal{F}_x \otimes \kappa(x)$ is nonzero.
By (2) this implies $x \in \text{Supp}(\mathcal{F})$ if and only
if $y \in \text{Supp}(f^*\mathcal{F})$, which is the content of
assertion (3).
\end{proof}

\end{document}
